\section{The vineyard and its fruit}
God's vineyard has been cared for, yet its fruit is a disappointment. Isaiah gives that disappointment a concrete name: injustice. Jesus brings the vineyard image into a confrontation with stewards who try to possess what was entrusted to them. Between the accusation and the parable, the assembly speaks as the damaged vine, asking God to return and give it life. The Sunday therefore brings a difficult conjunction into prayer: the people who owe God fruit also need God to make them live.

Two readings of this conjunction unfold here. The first follows the priority of gift: God gives, chooses, nourishes, and therefore calls forth a response that takes the form of justice and charity. Augustine's account of Christ's choosing in John and Chrysostom's account of the vineyard's prepared abundance illuminate that order. The second listens to a wounded people's appeal for restoration. Augustine and Robert Bellarmine read the psalm's vine as perfected in Christ; Chrysostom's Eucharistic exposition presses that restoration toward the difficult unity of actual neighbors.

The readings meet in Christ, the rejected Son and cornerstone. Their emphases remain different. One asks how fruit becomes possible; the other asks what becomes of a people whose life is damaged. Both must face judgment, and neither can turn divine generosity into indifference to what people do. The petition after Communion brings each toward a changed life: the sacrament is received so that its gift may shape the receiver.

The apostolic reading adds its own voice. Philippians moves from petition and thanksgiving to disciplined attention and practiced virtue. Its counsel of peace belongs to Paul's continuing address to the Philippians. Isaiah and the Gospel form the Sunday's vineyard pairing; Philippians is part of the semi-continuous apostolic course. The common Missal prayers and antiphons serve this Sunday in all three cycles. Their contribution here arises through their words and ritual places, without making every text say the same thing.\footnote{\work{General Introduction to the Lectionary}, nos. 66--68, 106--107; U.S. Lectionary 139.}

The two Communion antiphons give alternative conclusions to this movement. Lamentations places the communicant among those who seek and wait for the good Lord. The alternative based on First Corinthians makes the many and the one body explicit. Waiting and incorporation are compatible Christian realities, but these are alternatives, not two successive antiphons in one prescribed sequence.
